%%% FIELD proof-width-identity
Fix \(S\in\mathcal P_k\), a table \(y\) compatible with the selected arm--cell means, a future policy \(\pi\in\mathcal F\), and a cell \(x\in\mathcal X\).  Write
\[
a=\pi(x),\qquad t=m_\pi.
\]
We first bound the conditional interval diameter in this cell.  If
\(T_{xa}(S)=\varnothing\), the range restriction \([0,1]\) in the envelope
definition gives
\[
u_{xa}(t;y)-\ell_{xa}(t;y)\leq 1=r_{\pi x}(S).                    \tag{1}
\]
Suppose next that \(a=a_0\) and \(T_{x a_0}(S)\neq\varnothing\).  By
neutral invariance and compatibility, all matching neutral anchors have one
common value \(y_0\in[0,1]\).  The neutral case of the envelope definition
uses Lipschitz constant zero, and therefore
\[
\ell_{x a_0}(t;y)=\max\{0,y_0\}=y_0,
\qquad
u_{x a_0}(t;y)=\min\{1,y_0\}=y_0.
\]
Consequently,
\[
u_{x a_0}(t;y)-\ell_{x a_0}(t;y)=0=r_{\pi x}(S).                 \tag{2}
\]

It remains to consider \(a\neq a_0\).  For any
\(c\in T_{xa}(S)\), if \(m_c\leq t\), the two envelope formulas imply
\[
\ell_{xa}(t;y)\geq y_{cx},
\qquad
u_{xa}(t;y)\leq y_{cx}+\Gamma(t-m_c).
\]
If \(m_c\geq t\), the same formulas imply
\[
\ell_{xa}(t;y)\geq y_{cx}-\Gamma(m_c-t),
\qquad
u_{xa}(t;y)\leq y_{cx}.
\]
In either case, also using that both endpoints lie in \([0,1]\),
\[
u_{xa}(t;y)-\ell_{xa}(t;y)
\leq \min\{1,\Gamma|t-m_c|\}=d_{\pi xc}.                         \tag{3}
\]
Action-mismatching candidates have cost one.  Taking the minimum of (3)
over matching anchors and including the dummy bound one, then combining
with (1)--(2), proves
\[
u_{x,\pi(x)}(m_\pi;y)-\ell_{x,\pi(x)}(m_\pi;y)
\leq r_{\pi x}(S).                                                \tag{4}
\]
The weights satisfy \(q_x=N_x/N>0\).  Summing (4), using the definitions of
\(L_S\), \(U_S\), and \(R_\pi\), and maximizing over the finite nonempty
set \(\mathcal F\) yields, for every compatible \(y\),
\[
\max_{\pi\in\mathcal F}
\{U_S(\pi;y)-L_S(\pi;y)\}
\leq \max_{\pi\in\mathcal F}R_\pi(S).
\]
Taking the supremum over \(Y\in\mathcal M_N\) in the definition of
\(\overline W(S)\) gives
\[
\overline W(S)\leq\max_{\pi\in\mathcal F}R_\pi(S).              \tag{5}
\]

We prove the reverse inequality by an explicit pair of fixed arrays.
Finiteness and nonemptiness of \(\mathcal F\) give a maximizer
\(\pi^\star\) of \(R_\pi(S)\).  For each cell \(x\), put
\[
a_x=\pi^\star(x),\qquad t=m_{\pi^\star},
\qquad w_x=r_{\pi^\star x}(S).
\]
If \(a_x\neq a_0\) and \(T_{x a_x}(S)=\varnothing\), let
\(g_x^-(s)=0\) and \(g_x^+(s)=1\).  If
\(a_x\neq a_0\), \(T_{x a_x}(S)\neq\varnothing\), and \(w_x=0\),
let both functions be identically zero.  In the remaining nonneutral case,
\(T_{x a_x}(S)\neq\varnothing\) and \(w_x>0\).  This forces
\(\Gamma>0\), and we define
\[
g_x^-(s)=\min\{w_x,\max\{0,\Gamma(s-t)\}\},
\qquad
g_x^+(s)=\min\{w_x,\max\{0,\Gamma(s-t)+w_x\}\}.                 \tag{6}
\]
The maps in (6) are compositions of a nondecreasing affine map with
clipping to \([0,w_x]\).  Hence they are bounded, nondecreasing, and
\(\Gamma\)-Lipschitz.  By the definition of \(w_x\), every matching
anchor \(c\in T_{x a_x}(S)\) satisfies
\[
w_x\leq \min\{1,\Gamma|m_c-t|\},
\qquad\text{so}\qquad \Gamma|m_c-t|\geq w_x.                    \tag{7}
\]
If \(m_c<t\), equations (6)--(7) give
\(g_x^-(m_c)=g_x^+(m_c)=0\); if \(m_c>t\), they give
\(g_x^-(m_c)=g_x^+(m_c)=w_x\).  The equality \(m_c=t\) would imply
\(w_x=0\), contrary to this case.  Thus the two functions agree at every
matching selected anchor, while
\[
g_x^-(t)=0,
\qquad
g_x^+(t)=w_x.                                                       \tag{8}
\]
For \(a_x=a_0\), take the identical zero constants if a matching neutral
anchor exists and take the two constants zero and one otherwise.  Neutral
invariance holds in both cases, the functions agree at all matching selected
anchors, and their target difference is again \(w_x\).

For every cell use the preceding pair for action \(a_x\), and use the
identical zero function for every other action.  Define homogeneous
fixed-population arrays by assigning these cell functions to every unit in
the cell.  Every nonneutral trajectory is bounded, nondecreasing, and
\(\Gamma\)-Lipschitz, and every neutral trajectory is constant.  Therefore
the two arrays lie in \(\mathcal M_N^{\mathrm{CO},\Gamma}\), and the bridge
lemma also places them in \(\mathcal M_N\).  For a selected candidate and
cell, either its action is \(a_x\), in which case the two schedules agree by
(7)--(8), or its action is different, in which case both outcomes are zero.
Thus the arrays produce exactly the same selected table \(y^\star\); in
fact, they produce the same realized outcome for every unit under every
assignment to the selected arms.  Their causal values at \(\pi^\star\)
satisfy
\[
V_{\mu^+}(\pi^\star)-V_{\mu^-}(\pi^\star)
=\sum_{x\in\mathcal X}q_xw_x
=R_{\pi^\star}(S).                                                   \tag{9}
\]
The sharp-common-endpoints theorem places both feasible values in the same
identified interval conditional on \(y^\star\).  Hence its width is at least
the separation in (9).  Since \(\pi^\star\) maximizes \(R_\pi(S)\), this
and (5) give
\[
\overline W(S)
=\max_{\pi\in\mathcal F}R_\pi(S)
=\max_{\pi\in\mathcal F}\sum_{x\in\mathcal X}q_x
\min\bigl(\{1\}\cup\{d_{\pi xc}:c\in S\}\bigr).                \tag{10}
\]
The common table \(y^\star\) is generated by each member of the constructed
pair, so the supremum in \(\overline W(S)\) is attained by two
observationally indistinguishable fixed arrays.  Because those arrays lie in
\(\mathcal M_N^{\mathrm{CO},\Gamma}\), while the bridge lemma says that
this subclass and \(\mathcal M_N\) generate the same mean surfaces, the
same exact width is obtained on the individual-level Lipschitz subclass.

Finally, \(q_x=N_x/N\), every score, and \(\Gamma\) are rational.  The
costs \(d_{\pi xc}\) are therefore rational, and (10) uses only minima,
sums, and maxima over finite sets.  Thus \(\overline W(S)\) is rational for
each \(S\).  The feasible family \(\mathcal P_k\) is finite and nonempty
because it contains \(\varnothing\).  Its minimum \(D^\star\) is therefore
rational and attained.

%%% FIELD def-finite-score-gamma-saturation
Let
\[
\mathcal G
=\left\{|m_\pi-m_c|:\ \pi\in\mathcal F,\ c\in\mathcal C,
\ \text{and for some }x\in\mathcal X,
\ c(x)=\pi(x)\neq a_0,\ |m_\pi-m_c|>0\right\}.
\]
Define the nonnegative rational quantities
\[
\Delta=
\begin{cases}
\min\mathcal G,&\mathcal G\neq\varnothing,\\
0,&\mathcal G=\varnothing,
\end{cases}
\qquad
\Gamma_{\mathrm{sat}}=
\begin{cases}
1/\Delta,&\Delta>0,\\
0,&\Delta=0.
\end{cases}                                                        \tag{11}
\]
Thus no division is made in the no-positive-gap case.  For each
\(\gamma\in\mathbb Q_{\geq0}\), a superscript or subscript \(\gamma\)
denotes the response class, envelopes, cost tensor, residuals, and
prospective loss obtained from their finite-\(\Gamma\) definitions after
replacing \(\Gamma\) by \(\gamma\).  In particular,
\[
d^{(\gamma)}_{\pi xc}=
\begin{cases}
1,&c(x)\neq\pi(x),\\
0,&c(x)=\pi(x)=a_0,\\
\min\{1,\gamma|m_\pi-m_c|\},&c(x)=\pi(x)\neq a_0,
\end{cases}
\]
\[
r^{(\gamma)}_{\pi x}(S)
=\min\bigl(\{1\}\cup\{d^{(\gamma)}_{\pi xc}:c\in S\}\bigr),
\qquad
R^{(\gamma)}_\pi(S)=\sum_{x\in\mathcal X}q_x
r^{(\gamma)}_{\pi x}(S),
\]
and \(\overline W_\gamma(S)\) is the causal prospective sharp-diameter
loss over the response class with Lipschitz constant \(\gamma\), while
\[
D^\star(\gamma)=\min_{S\in\mathcal P_k}\overline W_\gamma(S).
\]

%%% FIELD stmt-finite-score-gamma-saturation
Let \(\Delta\) and \(\Gamma_{\mathrm{sat}}\) be defined by the finite
future-policy/candidate score gaps.  For every
\(\gamma\in\mathbb Q_{\geq0}\) satisfying
\(\gamma\geq\Gamma_{\mathrm{sat}}\),
\[
d^{(\gamma)}_{\pi xc}=d^{\mathrm{CO}}_{\pi xc}
\quad\text{for every }(\pi,x,c)\in
\mathcal F\times\mathcal X\times\mathcal C.                       \tag{12}
\]
Consequently, for every \(S\in\mathcal P_k\),
\[
\overline W_\gamma(S)=\overline W^{\mathrm{CO}}(S),                \tag{13}
\]
and therefore
\[
D^\star(\gamma)=\min_{S\in\mathcal P_k}\overline W^{\mathrm{CO}}(S),
\qquad
\operatorname*{arg\,min}_{S\in\mathcal P_k}\overline W_\gamma(S)
=\operatorname*{arg\,min}_{S\in\mathcal P_k}
\overline W^{\mathrm{CO}}(S).                                      \tag{14}
\]
Substitution of the binary tensor \(d^{\mathrm{CO}}\) for
\(d^{(\gamma)}\) therefore transfers the exact mixed-integer formulation,
the independently checked rational certificate workflow, and its accepted
certificates without changing their panel objective values.  The map
\(\gamma\mapsto D^\star(\gamma)\) on nonnegative rational \(\gamma\) is
the restriction of a continuous nondecreasing finite piecewise-linear
function on \(\mathbb R_{\geq0}\); its finite arrangement is induced by
the positive score-gap saturation knots and the crossings of the resulting
affine panel--policy pieces.  It is constant for
\(\gamma\geq\Gamma_{\mathrm{sat}}\).  If there is no positive same-action
nonneutral score gap, then \(\Delta=\Gamma_{\mathrm{sat}}=0\), and
(12)--(14) hold for every \(\gamma\geq0\).  These equalities unify the
finite-library design objectives; they do not assert equality of the
underlying aggregate and individual-level response classes.

%%% FIELD proof-finite-score-gamma-saturation
Let \(\mathcal G\), \(\Delta\), and \(\Gamma_{\mathrm{sat}}\) be as in
the finite-score saturation definition, and fix
\(\gamma\in\mathbb Q_{\geq0}\) with
\(\gamma\geq\Gamma_{\mathrm{sat}}\).  Fix a coordinate
\((\pi,x,c)\).  If \(c(x)\neq\pi(x)\), both
\(d^{(\gamma)}_{\pi xc}\) and \(d^{\mathrm{CO}}_{\pi xc}\) equal one
by their definitions.  If \(c(x)=\pi(x)=a_0\), both equal zero.  Suppose
therefore that
\[
c(x)=\pi(x)\neq a_0,
\qquad \delta=|m_\pi-m_c|.
\]
When \(\delta=0\), both costs again equal zero.  When \(\delta>0\), the
definition of \(\mathcal G\) gives \(\delta\in\mathcal G\); hence
\(\mathcal G\neq\varnothing\), \(\Delta>0\), and
\(\delta\geq\Delta\).  The threshold condition now gives
\[
\gamma\delta\geq \Gamma_{\mathrm{sat}}\Delta
=\Delta^{-1}\Delta=1.
\]
Thus \(d^{(\gamma)}_{\pi xc}=\min\{1,\gamma\delta\}=1\), which is also
the no-Lipschitz binary cost because the scores differ.  This exhausts the
cases and proves (12).  If \(\mathcal G=\varnothing\), the case
\(\delta>0\) cannot occur for a same-action nonneutral coordinate;
therefore the same case analysis proves (12) for every \(\gamma\geq0\)
while (11) sets \(\Gamma_{\mathrm{sat}}=0\) without division.

Taking minima over the same finite set of selected candidates in (12),
then multiplying by \(q_x>0\) and summing, gives
\[
r^{(\gamma)}_{\pi x}(S)=r^{\mathrm{CO}}_{\pi x}(S),
\qquad
R^{(\gamma)}_\pi(S)=R^{\mathrm{CO}}_\pi(S)                           \tag{15}
\]
for every \(S\in\mathcal P_k\), \(\pi\in\mathcal F\), and
\(x\in\mathcal X\).  The finite-\(\gamma\) width identity and the
no-Lipschitz Chen--Oberst sharpness theorem identify the two causal losses
with the respective computable maxima.  Equation (15) therefore yields
\[
\overline W_\gamma(S)
=\max_{\pi\in\mathcal F}R^{(\gamma)}_\pi(S)
=\max_{\pi\in\mathcal F}R^{\mathrm{CO}}_\pi(S)
=\overline W^{\mathrm{CO}}(S),                                      \tag{16}
\]
which proves (13).  Minimizing the pointwise equality (16) over the same
finite nonempty feasible family \(\mathcal P_k\) proves both the value and
full minimizer-set equalities in (14).

The panel mixed-integer program uses the cost tensor only through its
rational objective coefficients, and its equivalence theorem identifies
its fixed-selection value with the corresponding maximum of weighted
nearest-candidate residuals.  At saturation, (12) makes the
finite-\(\gamma\) and binary Chen--Oberst programs identical after literal
substitution.  Likewise, the certificate checker recomputes the incumbent
width and each optimistic descendant width; by (16), every such exact
rational comparison has the same value after substitution.  Hence a branch
tree is accepted in one formulation exactly when the substituted tree is
accepted in the other.  The certificate soundness-and-completeness theorem
turns checker acceptance into exact global optimality and guarantees a
finite accepted branch tree.  The exact-rational
generator-and-compression theorem supplies a terminating deterministic
generator and independently executed checker for the finite-\(\gamma\)
tensor.  Since their evaluated costs, branch decisions, incumbent values,
and acceptance tests are unchanged under (12) and (16), the entire
workflow transfers to the binary tensor, with the same optimal panels and
certificates.

It remains to prove the shape of the sensitivity path.  Temporarily allow
\(\gamma\in\mathbb R_{\geq0}\) in the explicit cost formula and define
the panel functions and real extension
\[
G_S(\gamma)=\max_{\pi\in\mathcal F}
\sum_{x\in\mathcal X}q_x r^{(\gamma)}_{\pi x}(S),
\qquad
\widetilde D(\gamma)=\min_{S\in\mathcal P_k}G_S(\gamma).            \tag{17}
\]
For fixed \((S,\pi,x)\), let
\[
M_{\pi x}(S)=\{c\in S:c(x)=\pi(x)\}.
\]
The score-cost definition gives the following exhaustive formula:
\[
r^{(\gamma)}_{\pi x}(S)=
\begin{cases}
0,&\pi(x)=a_0\text{ and }M_{\pi x}(S)\neq\varnothing,\\
1,&M_{\pi x}(S)=\varnothing,\\
\min\{1,\gamma\delta_{\pi x}(S)\},
&\pi(x)\neq a_0\text{ and }M_{\pi x}(S)\neq\varnothing,
\end{cases}                                                         \tag{18}
\]
where, in the last line, the finite nonempty set permits the definition
\[
\delta_{\pi x}(S)=\min_{c\in M_{\pi x}(S)}|m_\pi-m_c|.
\]
Every function in (18) is continuous, nondecreasing, and piecewise affine,
with at most one positive saturation knot, namely
\(1/\delta_{\pi x}(S)\) when \(\delta_{\pi x}(S)>0\).  Every positive
\(\delta_{\pi x}(S)\) is one of the finitely many score gaps in
\(\mathcal G\).  On each interval cut out by these reciprocal gaps, every
weighted policy sum is affine.  A finite maximum over policies and a finite
minimum over panels remain continuous and piecewise affine after refining
those intervals at the finitely many crossings of the affine pieces.
Therefore (17) is a continuous finite piecewise-linear function.  Since
the scores and weights are rational, all nondegenerate knots and crossings
are rational.

For \(0\leq\gamma_1\leq\gamma_2\), equation (18) gives
\(r^{(\gamma_1)}_{\pi x}(S)\leq
r^{(\gamma_2)}_{\pi x}(S)\).  Positive weighted sums and maxima preserve
this inequality, so
\[
G_S(\gamma_1)\leq G_S(\gamma_2)
\quad\text{for every }S.                                            \tag{19}
\]
Let \(S_2\) minimize the right-hand criterion at \(\gamma_2\), which
exists because \(\mathcal P_k\) is finite and nonempty.  Then
\[
\widetilde D(\gamma_1)
\leq G_{S_2}(\gamma_1)
\leq G_{S_2}(\gamma_2)
=\widetilde D(\gamma_2).                                            \tag{20}
\]
Thus the path is nondecreasing.  For rational \(\gamma\), the width
identity makes the causal optimum \(D^\star(\gamma)\) equal to (17), so
the latter is the asserted real piecewise-linear extension.  Finally, the
coordinatewise case analysis proving (12) uses only real inequalities and
therefore also applies to the real extension.  Thus (15), with
\(G_S(\gamma)\) in place of the causal loss, makes (17) equal to the same
binary optimum for every \(\gamma\geq\Gamma_{\mathrm{sat}}\); hence it is
constant from the threshold onward.  All these conclusions compare the exact design
functionals over the fixed finite libraries.  They use neither, and imply
neither, equality of the aggregate response class with the unrestricted
individual-level Chen--Oberst class.
